% 稳健性检验
\section{稳健性检验}

为确保基准结论的可靠性，本节从五个维度进行稳健性检验。表 \ref{tab:robustness} 汇总了六组回归的核心结果（R1 为基准 M6，用于对比）。

\begin{table}[htbp]\centering
\def\sym#1{\ifmmode^{#1}\else\(^{#1}\)\fi}
\caption{稳健性检验\label{tab:robustness}}
\begin{tabular}{l*{6}{c}}
\toprule
                    &\multicolumn{1}{c}{(1)}&\multicolumn{1}{c}{(2)}&\multicolumn{1}{c}{(3)}&\multicolumn{1}{c}{(4)}&\multicolumn{1}{c}{(5)}&\multicolumn{1}{c}{(6)}\\
                    &\multicolumn{1}{c}{基准}&\multicolumn{1}{c}{替换D}&\multicolumn{1}{c}{剔除直辖市}&\multicolumn{1}{c}{对数变换}&\multicolumn{1}{c}{剔除COVID}&\multicolumn{1}{c}{5\%缩尾}\\
\midrule
digital\_economy\_index&       0.277\sym{***}&                     &       0.142\sym{*}  &       0.252\sym{***}&       0.473\sym{***}&       0.242\sym{***}\\
                    &     (0.085)         &                     &     (0.070)         &     (0.077)         &     (0.136)         &     (0.063)         \\
\addlinespace
digital\_finance\_index&                     &       0.055         &                     &                     &                     &                     \\
                    &                     &     (0.037)         &                     &                     &                     &                     \\
\addlinespace
ln(人均GDP)         &                     &                     &                     &      -0.093         &                     &                     \\
                    &                     &                     &                     &     (0.064)         &                     &                     \\
\midrule
Observations        &         403         &         403         &         364         &         403         &         372         &         403         \\
\bottomrule
\multicolumn{7}{l}{\footnotesize 注：所有模型均包含控制变量、省份固定效应和年份固定效应，省份聚类标准误。}\\
\multicolumn{7}{l}{\footnotesize R1基准模型，R2替换核心解释变量为数字金融指数，R3剔除北京/天津/上海三个直辖市。}\\
\multicolumn{7}{l}{\footnotesize R4对数变换人均GDP和教育支出，R5剔除2020年，R6使用5\%缩尾替代1\%。}\\
\multicolumn{7}{l}{\footnotesize * p$<$0.10, ** p$<$0.05, *** p$<$0.01}\\
\end{tabular}
\end{table}


\subsection{替换核心解释变量}

将数字经济综合指数替换为北大数字普惠金融指数后（R2），核心系数从前述的 0.277 下降至 0.055，且在 10\% 水平上不显著（$p=0.153$）。这一"消失的效应"具有两个层面的含义。在方法层面，它表明本文的核心结论——数字经济对就业结构的影响——依赖于广义数字经济综合指数的测度，对单纯的数字金融指数不稳健。这提示，如果读者认同数字金融指数是数字经济的合理代理变量，则基准结论需要打折扣。在实质层面，它揭示了一个可能的经济学差异：广义数字经济（覆盖互联网基础设施、移动通信、软件服务业等多维度）通过改变生产方式、降低交易成本、扩展市场半径等渠道重塑就业结构，而数字金融的发展主要通过降低融资约束、促进创业活动来发挥作用——这两个渠道在推动劳动力从制造业向服务业转移上的力道可能不同。从政策优先序角度看，这一差异意味着，如果目标是推动就业结构转型，宽带提速和 5G 覆盖可能比推广移动支付和网络借贷有更强的直接效应。

\subsection{剔除直辖市}

将北京、天津和上海三个直辖市从样本中剔除后（R3），数字经济系数下降至 0.142，显著性降至边缘水平（$p=0.051$）。这一变化直接反映了三个直辖市在全样本中的杠杆效应：它们在数字经济和三产就业占比两个维度均处于样本顶端，为全样本回归提供了大量的识别变异。剔除后系数缩减约 49\%，标准误保持稳定。有两个互补的解读方式：（1）基准结果对直辖市观测有一定依赖，读者在将结论推广至"非直辖市省份"时应当审慎；（2）剔除直辖市后系数仍然为正且在 10\% 水平上显著，且样本量从 403 降至 364（下降仅 10\%），核心关系的正号在中短期内不需要直辖市来支撑。

\subsection{对数变换}

将人均 GDP 和教育支出进行对数变换以缓解潜在的异方差和非线性问题后（R4），数字经济系数为 0.252（$p<0.01$），与基准系数 0.277 在量级上高度一致。对数变换改变了控制变量的函数形式，但核心系数仅变动了约 9\%，表明基准结果对变量函数形式的选择不敏感。

\subsection{剔除 COVID-19 冲击年份}

2020 年 COVID-19 疫情的暴发导致全国范围内服务业大幅萎缩——餐饮、旅游、交通等第三产业部门首当其冲，三产就业占比出现异常性波动。与此同时，疫情防控催生了线上办公、远程教育、数字医疗等场景的加速渗透，数字经济指数可能不降反升。如果 2020 年的数据同时包含了"异常低的三产就业"和"异常高的数字经济指数"，则会对回归系数产生下行偏误。剔除 2020 年观测值后（R5），数字经济系数上升至 0.473（$p<0.01$），约为基准系数的 1.7 倍。这一结果证实了上述担忧：COVID-19 冲击确实压低了数字经济的就业结构效应估计。从中可以提取出一个有政策洞见的信息：数字经济促进就业结构转型的机制需要经济活动的正常运转作为前提——当疫情封锁使大量服务业企业无法开门营业时，即使线上替代活动增加，数字经济也无法有效转化为服务业就业的增长。

\subsection{替代缩尾水平}

将缩尾比例从基准的 1\%/99\% 调整为 5\%/95\% 后（R6），数字经济系数为 0.242（$p<0.01$），与基准系数的差异在统计上不显著。这表明基准结果不是由极端值的任意剔除决策驱动的。

\subsection{稳健性检验小结}

综合五项稳健性检验，可以得出以下判断：基准结论——数字经济发展对三产就业占比有显著正向效应——在替换缩尾水平、对数变换和剔除 COVID 年份后保持稳健。主要的脆弱点存在于两个检验：替换核心解释变量为数字金融指数后效应消失，以及剔除直辖市后显著性降至边缘水平。这两个脆弱点的经济含义已在对应段落中讨论，它们不构成对基准结论的否定，而是界定了结论的适用范围和条件：广义数字经济（非狭义数字金融）对非直辖市省份的就业结构效应值得进一步检验。
